\documentclass[11pt,a4paper,oneside]{book}
\usepackage[utf8]{inputenc}
\usepackage[russian]{babel}
\usepackage{geometry}
\geometry{left=2.5cm,right=2.5cm,top=2.5cm,bottom=2.5cm}
\usepackage{fancyhdr}
\usepackage{graphicx}
\usepackage{amsmath}
\usepackage{amssymb}
\usepackage{listings}
\usepackage{color}
\usepackage{xcolor}
\usepackage{hyperref}
\usepackage{booktabs}
\usepackage{array}
\usepackage{longtable}
\usepackage{multirow}
\usepackage{tikz}
\usepackage{verbatim}
\usepackage{framed}
\usepackage{setspace}
% Margins and spacing
\setstretch{1.05}
\setlength{\parindent}{0pt}
\setlength{\parskip}{4pt}
% Aggressive chapter/section spacing reduction
\usepackage{titlesec}
\titlespacing*{\part}{0pt}{0pt}{0pt}
\titlespacing*{\chapter}{0pt}{0pt}{0pt}
\titlespacing*{\section}{0pt}{0pt}{0pt}
\titlespacing*{\subsection}{0pt}{0pt}{0pt}
\titlespacing*{\subsubsection}{0pt}{0pt}{0pt}
% Code style
\definecolor{codegreen}{rgb}{0,0.6,0}
\definecolor{codegray}{rgb}{0.5,0.5,0.5}
\definecolor{codepurple}{rgb}{0.58,0,0.82}
\definecolor{backcolour}{rgb}{0.95,0.95,0.95}
\definecolor{darkblue}{rgb}{0,0,0.5}
\lstset{
  backgroundcolor=\color{backcolour},
  commentstyle=\color{codegreen},
  keywordstyle=\color{magenta},
  numberstyle=\tiny\color{codegray},
  stringstyle=\color{codepurple},
  basicstyle=\ttfamily\small,
  breakatwhitespace=false,
  breaklines=true,
  captionpos=b,
  keepspaces=true,
  numbers=left,
  numbersep=5pt,
  showspaces=false,
  showstringspaces=false,
  showtabs=false,
  tabsize=2,
  frame=single,
  rulecolor=\color{black},
  language=TypeScript,
  inputencoding=utf8,
  extendedchars=true
}
% Header and footer
\pagestyle{fancy}
\fancyhf{}
\fancyhead[L]{\textit{Stvor Messenger}}
\fancyhead[R]{\textit{Архитектура v1.0}}
\fancyfoot[C]{\thepage}
\renewcommand{\headrulewidth}{0.4pt}
\renewcommand{\footrulewidth}{0.4pt}
% Title page
\title{%
  \textbf{STVOR MESSENGER} \\
    {\large Полный архитектурный анализ} \\
    {\large Ilyazh-Web3E2E Protocol v0.9.0-beta}
}
\author{%
  Анализ кодовой базы \\
    Версия: 0.8.0 \\
  Статус: Beta
}
\date{Декабрь 2025}
% Table of contents depth
\setcounter{tocdepth}{2}
\begin{document}
\maketitle
\begin{center}
  \Large{\textbf{Полная архитектурная документация}} \\
    Post-Quantum E2E Encrypted Messenger \\
    123 KB документации • 3,136 строк • ~25,000 слов \\
  Анализ каждого файла и строки кода
\end{center}
\tableofcontents
\pagebreak
% PART I: FUNDAMENTALS
\vspace{-1cm}\part{Фундаментальные принципы}
\vspace{-0.5cm}\chapter{Введение в Stvor Messenger}
Stvor Messenger — это \textbf{пост-квантовый E2E-зашифрованный мессенджер}, построенный на протоколе \textbf{Ilyazh-Web3E2E v0.8}. Проект демонстрирует правильное проектирование безопасного мессенджера с учётом будущих квантовых угроз.
\vspace{-0.3cm}\section{Три столпа архитектуры}
\vspace{-0.2cm}\subsection{1. Post-Quantum Cryptography (PQC)}
\textbf{Цель:} Защита от атак «Harvest Now, Decrypt Later» (HNDL)
\textbf{Проблема:} Если злоумышленник записывает зашифрованные сообщения сегодня, он может расшифровать их в будущем с помощью квантового компьютера.
\textbf{Решение:} Гибридная криптография (классическая + постквантовая одновременно):
\begin{lstlisting}[language=]
NOT:  ECDH OR PQ_CRYPTO
YES:  ECDH AND PQ_CRYPTO (одновременно)
shared_secret = HKDF(ecdh_secret || mlkem_secret)
\end{lstlisting}
Если один механизм скомпрометирован, другой остаётся безопасным.
\vspace{-0.2cm}\subsection{2. Client-Side Encryption}
\begin{itemize}
  \item Приватные ключи \textbf{НИКОГДА} не покидают браузер
  \item Сервер реле хранит только \textbf{общедоступные ключи} и \textbf{зашифрованные сообщения}
  \item Дешифрование происходит полностью на клиенте
  \item Relay = stateless message broker
\end{itemize}
\vspace{-0.2cm}\subsection{3. Defense-in-Depth}
Множественные независимые слои защиты:
\begin{enumerate}
  \item \textbf{Transport Layer} — TLS 1.3
  \item \textbf{Relay Integrity} — Identity pinning (EREBUS mitigation)
  \item \textbf{Message Encryption} — AEAD with AAD
  \item \textbf{Traffic Analysis} — Message padding, timing jitter
  \item \textbf{User Privacy} — Toggles for typing, read receipts
  \item \textbf{PQ Downgrade} — Mandatory PQ detection
\end{enumerate}
\vspace{-0.5cm}\chapter{Криптографический протокол}
\vspace{-0.3cm}\section{Общая архитектура Ilyazh-Web3E2E}
\vspace{-0.2cm}\subsection{Layer 1: Handshake (Authenticated Key Exchange)}
Гибридный обмен ключами с двойными подписями:
\begin{lstlisting}[language=]
Client A                          Client B
1. Fetch B's prekeys from relay
   → identity: Ed25519 + ML-DSA pub keys
   → signed prekey: X25519 + ML-KEM pub keys
2. A generates ephemeral keys
   → X25519 ephemeral keypair
   → ML-KEM-768 ephemeral keypair
3. A creates handshake
   → ECDH(A.eph_x25519, B.id_x25519)
   → KEM_encap(B.mlkem_pubkey) → shared_secret, ct
   → Sign with Ed25519 + ML-DSA-65
4. A sends handshake to relay (POST /api/message/:chatId)
   → messageType: 'handshake'
   → Relay stores as blob with blobRef = SHA256(data)
5. B receives handshake (via WebSocket or polling)
   → ECDH(B.id_x25519_secret, A.eph_x25519_pub)
   → KEM_decap(B.mlkem_secret, A.mlkem_ct)
   → Verify Ed25519 signature
   → Verify ML-DSA signature
   → Compute session ID = SHA256(transcript)
6. B sends response handshake
   → Same structure, no ephemeralMLKEM/ciphertext fields
7. A receives response
   → Verifies all B's signatures
   → Sets pqMandatory = true (both support PQ)
   → Session ready
✅ E2E encrypted session established
\end{lstlisting}
\vspace{-0.2cm}\subsection{Layer 2: Double Ratchet}
\subsubsection{Ratchet State}
\begin{lstlisting}[language=TypeScript]
interface RatchetState {
  sessionId: Uint8Array              // 32 bytes, unique per session
  rootKey: Uint8Array                // 64 bytes, from ECDH + KEM
  sendChainKey: Uint8Array           // 64 bytes, evolves per message
  recvChainKey: Uint8Array           // 64 bytes, evolves per message
  sendRatchetId: bigint              // 8 bytes, epoch counter
  recvRatchetId: bigint              // 8 bytes, epoch counter
  sendCounter: number                // 4 bytes, per-epoch counter
  recvCounter: number                // 4 bytes, per-epoch counter
  epochStartTime: number             // Timestamp when epoch began
  sessionStartTime: number            // Timestamp when session started
  totalMessages: number              // Total messages in session
  pqMandatory: boolean               // Downgrade protection
}
\end{lstlisting}
\subsubsection{Mandated Re-encapsulation Cadence}
Обязательный DH ratchet каждые:
\begin{itemize}
  \item \textbf{2\textsuperscript{20}} (1M) сообщений в эпохе, ИЛИ
  \item \textbf{24 часа} с начала эпохи, ИЛИ
  \item \textbf{2\textsuperscript{32}} (4B) сообщений за сессию, ИЛИ
  \item \textbf{7 дней} с начала сессии
\end{itemize}
\vspace{-0.2cm}\subsection{Layer 3: Per-Message Encryption (AEAD)}
\subsubsection{Для каждого сообщения:}
\begin{enumerate}
  \item \textbf{Проверка переключения (rekeying check)}
  \begin{lstlisting}[language=TypeScript]
needsRekey(state) → {
  required: boolean,
  reason?: 'message_limit' | 'time_limit' | 'session_cap_messages' | 'session_cap_time'
}
  \end{lstlisting}
  \item \textbf{Вычисление ключа сообщения}
  \begin{lstlisting}[language=]
mk = HKDF-SHA256(sendChainKey, label="mk", 32)
  \end{lstlisting}
  \item \textbf{Построение AAD (Additional Authenticated Data)}
  \begin{lstlisting}[language=]
AAD = Version(1) || SuiteID(8) || SessionID(32) ||
      Sequence(8) || Epoch(8) || Flags(1)
= 58 байт total
  \end{lstlisting}
  \item \textbf{Конструирование nonce}
  \begin{lstlisting}[language=]
nonce = ratchetId || counter
= 8 bytes || 4 bytes = 12 bytes total (GCM требует 12)
  \end{lstlisting}
  \item \textbf{Padding сообщения (defense-in-depth)}
  \begin{lstlisting}[language=TypeScript]
padded = padMessage(plaintext, {
  blockSize: 256 | 512 | 1024,  // Align to power-of-2
  strategy: 'random' | 'pkcs7'
})
  \end{lstlisting}
  \item \textbf{Шифрование AEAD}
  \begin{lstlisting}[language=]
ciphertext = AES-256-GCM(mk, nonce, padded_plaintext, aad)
Output: ciphertext + 16-byte auth tag
  \end{lstlisting}
  \item \textbf{Отправка на relay}
  \begin{lstlisting}[language=TypeScript]
blobRef = SHA256(ciphertext_blob)
POST /api/message/:chatId → {
  blobRef,
  encryptedBlob,
  timestamp
}
Relay response: { sequence, timestamp }
  \end{lstlisting}
\end{enumerate}
\vspace{-0.2cm}\subsection{Размеры компонентов криптопротокола}
\begin{table}[h]
\centering
\begin{tabular}{|l|r|l|l|}
\hline
\textbf{Компонент} & \textbf{Размер} & \textbf{Цель} & \textbf{Статус} \\
\hline
X25519 Pub/Sec & 32/32 & ECDH & Классический \\
\hline
ML-KEM-768 Pub & 1184 & PQ инкапсуляция & FIPS 203 \\
\hline
ML-KEM-768 Ct & 1088 & PQ общий секрет & FIPS 203 \\
\hline
Ed25519 Sig & 64 & Подпись & Классический \\
\hline
ML-DSA-65 Sig & 3309 & PQ подпись & FIPS 204 \\
\hline
Session ID & 32 & ID сессии (AAD) & Производный \\
\hline
Chain Key & 64 & Состояние рэтчета & Производный \\
\hline
\end{tabular}
\end{table}
% PART II: ARCHITECTURE
\vspace{-1cm}\part{Архитектура системы}
\vspace{-0.5cm}\chapter{Общая схема}
\vspace{-0.3cm}\section{Трёхуровневая архитектура}
\begin{verbatim}
┌──────────────────────────────┐
│   FRONTEND (Next.js 16)      │
│   (apps/web, 6,963 строк)   │
├──────────────────────────────┤
│ • React Components           │
│ • Clerk Auth                 │
│ • IndexedDB Keystore         │
│ • WebSocket Client           │
│ • Crypto Initialization      │
└──────────────┬───────────────┘
               │ HTTPS + WSS (TLS 1.3)
┌──────────────▼───────────────┐
│  RELAY SERVER (Fastify)      │
│  (apps/relay, 5,037 строк)  │
├──────────────────────────────┤
│ • Stateless message broker   │
│ • WebSocket manager          │
│ • Rate limiting              │
│ • API routes                 │
│ • Storage adapter (pluggable)│
└──────────────┬───────────────┘
               │ SQL / In-Memory
┌──────────────▼───────────────┐
│    STORAGE LAYER             │
├──────────────────────────────┤
│ PostgreSQL (production)      │
│ In-Memory (development)      │
└──────────────────────────────┘
\end{verbatim}
\vspace{-0.3cm}\section{Монорепозиторий (pnpm + Turbo)}
\begin{verbatim}
ilyazh-messenger/
├── apps/
│   ├── relay/           ← Backend Fastify сервер
│   │   ├── src/
│   │   │   ├── index.ts (2,211 LOC)
│   │   │   ├── websocket.ts (612 LOC)
│   │   │   ├── storage.ts (335 LOC)
│   │   │   └── storage/
│   │   │       ├── interfaces.ts (6,825 LOC)
│   │   │       ├── memory-adapter.ts
│   │   │       └── postgres-adapter.ts
│   │   └── Dockerfile
│   │
│   └── web/             ← Frontend Next.js
│       ├── app/
│       │   ├── page.tsx
│       │   ├── layout.tsx
│       │   ├── (dashboard)/ (protected routes)
│       │   └── api/ (route handlers)
│       └── lib/ (6,963 LOC)
│           ├── crypto/init.ts (9,698 LOC)
│           ├── identity.ts (672 LOC)
│           ├── keystore.ts (685 LOC)
│           └── ... (25+ more files)
│
├── packages/
│   └── crypto/          ← Shared crypto library
│       ├── src/ (5,995 LOC)
│       │   ├── constants.ts (57 LOC)
│       │   ├── primitives.ts (703 LOC)
│       │   ├── handshake.ts (976 LOC)
│       │   ├── ratchet.ts (389 LOC)
│       │   ├── defense-in-depth.ts (838 LOC)
│       │   └── ... (10+ more files)
│       └── dist/ (compiled + .d.ts)
│
└── turbo.json           ← Build orchestration
\end{verbatim}
\vspace{-0.5cm}\chapter{Backend: Relay Server}
\vspace{-0.3cm}\section{Архитектура реле-сервера}
Relay Server является \textbf{stateless message broker}. Может масштабироваться горизонтально на 100+ инстансов.
\vspace{-0.2cm}\subsection{Основные компоненты (apps/relay/src/index.ts)}
\begin{enumerate}
  \item \textbf{Plugin Registration}
  \begin{itemize}
    \item CORS (configurable allowed origins)
    \item JWT (token validation)
    \item Multipart (file uploads)
    \item Rate Limiting (sliding window: 100 req/min)
    \item WebSocket (real-time)
  \end{itemize}
  \item \textbf{Storage Adapter (Pluggable)}
  \begin{itemize}
    \item Interface: IStorageAdapter
    \item Memory adapter (development)
    \item PostgreSQL adapter (production)
  \end{itemize}
  \item \textbf{API Routes}
  \begin{itemize}
    \item POST /api/register
    \item GET /api/prekeys/:userId
    \item POST /api/message/:chatId
    \item GET /api/sync
    \item GET /api/feed
    \item WebSocket /ws
  \end{itemize}
  \item \textbf{Business Logic}
  \begin{itemize}
    \item User management (registration, lookup)
    \item Prekey management (storage, rotation, consumption)
    \item Message storage (content-addressed blobs)
    \item Sync cursors (per-user polling state)
    \item Rate limiting (per-user sliding window)
  \end{itemize}
\end{enumerate}
\vspace{-0.3cm}\section{Storage Abstraction Layer}
Repository Pattern с 6 интерфейсами (apps/relay/src/storage/interfaces.ts):
\vspace{-0.2cm}\subsection{IUserRepository}
\begin{lstlisting}[language=TypeScript]
interface UserIdentity {
  userId: string
  username: string
  identityEd25519: string    // base64
  identityMLDSA: string      // base64
  registeredAt: number
}
Methods:
- createUser(user: UserIdentity): Promise<boolean>
- getUserById(userId: string): Promise<UserIdentity | null>
- getUserByUsername(username: string): Promise<UserIdentity | null>
- userExists(userId: string): Promise<boolean>
- getTotalUserCount(): Promise<number>
\end{lstlisting}
\vspace{-0.2cm}\subsection{IPrekeyRepository}
\begin{lstlisting}[language=TypeScript]
interface PrekeyBundle {
  userId: string
  bundleId: string           // UUID
  x25519Ephemeral: string    // base64
  mlkemPublicKey: string     // base64
  ed25519Signature: string   // base64
  mldsaSignature: string     // base64
  timestamp: number
  isOneTime: boolean
  consumed: boolean
  consumedAt?: number
  consumedBy?: string        // userId
}
Methods:
- storePrekeyBundle(bundle: PrekeyBundle): Promise<boolean>
- getSignedPrekey(userId: string): Promise<PrekeyBundle | null>
- consumeOneTimePrekey(userId, consumedBy): Promise<PrekeyBundle | null>
- listPrekeyBundles(userId: string): Promise<PrekeyBundle[]>
- deleteOldPrekeys(userId, olderThan): Promise<number>
- countAvailableOneTimePrekeys(userId): Promise<number>
\end{lstlisting}
\vspace{-0.2cm}\subsection{IMessageRepository (Content-Addressed)}
\begin{lstlisting}[language=TypeScript]
interface MessageBlob {
  chatId: string
  blobRef: string            // SHA-256 hash
  encryptedBlob: Buffer      // ENCRYPTED bytes
  createdAt: number
  senderId: string
  sequence: number           // Monotonic per chat
  messageType?: 'handshake' | 'message'
}
Methods:
- storeBlob(blob: MessageBlob): Promise<boolean>
- getBlob(blobRef: string): Promise<MessageBlob | null>
- listBlobsByChatId(chatId, limit): Promise<MessageBlob[]>
- deleteBlobsByChatId(chatId): Promise<number>
- deleteOldBlobs(olderThan: number): Promise<number>
\end{lstlisting}
\vspace{-0.2cm}\subsection{ISyncRepository (Polling Cursors)}
\begin{lstlisting}[language=TypeScript]
interface SyncCursor {
  userId: string
  chatId: string
  lastSeenSequence: number
  updatedAt: number
}
Methods:
- setSyncCursor(cursor: SyncCursor): Promise<void>
- getSyncCursor(userId, chatId): Promise<SyncCursor | null>
- getUnseenMessageCount(userId, chatId): Promise<number>
\end{lstlisting}
\vspace{-0.2cm}\subsection{IRateLimitRepository}
\begin{lstlisting}[language=TypeScript]
Methods:
- checkLimit(userId, window, limit): Promise<boolean>
- recordRequest(userId): Promise<void>
- resetLimit(userId): Promise<void>
\end{lstlisting}
\vspace{-0.2cm}\subsection{IPostRepository (News Feed)}
\begin{lstlisting}[language=TypeScript]
interface Post {
  postId: string             // UUID
  authorId: string
  authorUsername: string
  content: string
  imageUrl?: string
  createdAt: number
  likesCount: number
  commentsCount: number
  sharesCount: number
}
Methods:
- createPost(post: Post): Promise<boolean>
- getPostById(postId): Promise<Post | null>
- listPosts(limit, offset): Promise<Post[]>
- updatePost(postId, updates): Promise<boolean>
- deletePost(postId): Promise<boolean>
\end{lstlisting}
\vspace{-0.3cm}\section{WebSocket Manager (612 LOC)}
Real-time messaging через WebSocket:
\begin{lstlisting}[language=TypeScript]
class WebSocketManager {
  private clients: Map<username, ConnectedClient>
  private chatSubscriptions: Map<chatId, Set<username>>
  private typingIndicators: Map<chatId, Map<username, timeout>>
  registerClient(username: string, ws: WebSocket, chatIds: string[]): void
  unregisterClient(username: string): void
  subscribeToChat(username: string, chatId: string): void
  unsubscribeFromChat(username: string, chatId: string): void
  broadcastToChat(chatId: string, message: WSMessage): void
  sendToClient(username: string, message: WSMessage): void
  sendTypingIndicator(chatId, username): void
  sendReadReceipt(chatId, username, sequence): void
  updatePresence(username, status): void
}
\end{lstlisting}
\vspace{-0.3cm}\section{Beta Limits}
\begin{lstlisting}[language=TypeScript]
const MAX_USERS = 100                              // Beta test limit
const MAX_MESSAGES_PER_CHAT = 1000                 // Per chat
const MAX_MESSAGE_SIZE = 1 * 1024 * 1024          // 1 MB
const MESSAGE_TTL_MS = 7 * 24 * 60 * 60 * 1000   // 7 days
\end{lstlisting}
\vspace{-0.5cm}\chapter{Frontend: Web Client}
\vspace{-0.3cm}\section{Архитектура приложения (Next.js 16)}
\vspace{-0.2cm}\subsection{Три основных слоя}
\begin{enumerate}
  \item \textbf{Authentication Layer}
  \begin{itemize}
    \item Clerk Provider (OAuth, 2FA, sessions)
    \item Middleware (CSP, auth checks)
    \item Protected routes ((dashboard) group)
  \end{itemize}
  \item \textbf{Crypto Layer}
  \begin{itemize}
    \item CryptoInitializer (PQ + Classical)
    \item Keystore (IndexedDB persistence)
    \item Identity (E2E keypairs)
    \item Prekeys (bundle generation)
    \item Ratchet State (session management)
  \end{itemize}
  \item \textbf{Communication Layer}
  \begin{itemize}
    \item WebSocket Client (real-time)
    \item REST Client (message sync)
    \item Relay Client (user lookup, prekeys)
    \item Message Encryption/Decryption
  \end{itemize}
\end{enumerate}
\vspace{-0.3cm}\section{Lib Модули (31 файл, 6,963 LOC)}
\vspace{-0.2cm}\subsection{Криптография \& Управление ключами}
\begin{table}[h]
\centering
\begin{tabular}{|l|r|l|}
\hline
\textbf{Файл} & \textbf{LOC} & \textbf{Функция} \\
\hline
crypto/init.ts & 9,698 & Инициализация crypto системы \\
\hline
identity.ts & 672 & Управление E2E идентичностью \\
\hline
keystore.ts & 685 & IndexedDB persistence \\
\hline
secure-keystore.ts & 601 & Encrypted key storage \\
\hline
prekeys.ts & 505 & Prekey bundle generation \\
\hline
ratchet-refresh.ts & 496 & State recovery \& management \\
\hline
\end{tabular}
\end{table}
\vspace{-0.2cm}\subsection{Безопасность \& Сессии}
\begin{table}[h]
\centering
\begin{tabular}{|l|r|l|}
\hline
\textbf{Файл} & \textbf{LOC} & \textbf{Функция} \\
\hline
session-security.ts & 244 & Validation \& health checks \\
\hline
session-serializer.ts & 155 & State persistence \\
\hline
secure-relay-token.ts & 275 & Auth token management \\
\hline
relay-identity.ts & 151 & Identity verification \\
\hline
defense-integration.ts & 523 & Defense-in-depth \\
\hline
production-guard.ts & 98 & Safety checks \\
\hline
\end{tabular}
\end{table}
\vspace{-0.2cm}\subsection{Коммуникация \& Синхронизация}
\begin{table}[h]
\centering
\begin{tabular}{|l|r|l|}
\hline
\textbf{Файл} & \textbf{LOC} & \textbf{Функция} \\
\hline
websocket-client.ts & 389 & WebSocket connection \\
\hline
use-chat-realtime.ts & 216 & React hook for real-time \\
\hline
message-store.ts & 246 & Client-side cache \\
\hline
message-entry-validator.ts & 294 & Validation rules \\
\hline
relay-register-backoff.ts & 235 & Exponential backoff \\
\hline
\end{tabular}
\end{table}
\vspace{-0.3cm}\section{API Routes}
\vspace{-0.2cm}\subsection{User Management}
\begin{lstlisting}[language=]
POST /api/register
  Body: { username, identityEd25519, identityMLDSA }
  Response: { userId, username }
GET /api/prekeys/:userId
  Query: { includeOneTime: boolean }
  Response: { signedPrekey, oneTimePrekey }
\end{lstlisting}
\vspace{-0.2cm}\subsection{Messages}
\begin{lstlisting}[language=]
POST /api/message/:chatId
  Body: { messageType, blobRef, encryptedBlob }
  Response: { sequence, timestamp }
GET /api/sync?chatId=X&lastSeq=N
  Response: [{ sequence, blobRef, sender, timestamp }]
\end{lstlisting}
\vspace{-0.2cm}\subsection{WebSocket}
\begin{lstlisting}[language=]
Message types:
  { type: 'message', chatId, content }
  { type: 'typing', chatId, username }
  { type: 'read', chatId, sequence }
  { type: 'presence', status }
  { type: 'ping' | 'pong' }
\end{lstlisting}
% PART III: SECURITY
\vspace{-1cm}\part{Безопасность}
\vspace{-0.5cm}\chapter{Defense-in-Depth}
\vspace{-0.3cm}\section{Четыре независимых слоя защиты}
\vspace{-0.2cm}\subsection{Layer 1: Relay Network Integrity (EREBUS)}
\textbf{Проблема:} MITM может подменить relay сервер
\textbf{Решение:} Application-layer identity pinning
\begin{verbatim}
Protocol:
1. Client generates nonce (32 bytes, random)
2. Client → Relay: { action: "relay_challenge", nonce }
3. Relay → Client: { action: "relay_response", signature, pubkey }
4. Client verifies:
   - SHA256(pubkey) == expectedIdentityKeyHash ✓
   - Ed25519.verify(signature, nonce, pubkey) ✓
Result: Client is certain relay is authentic (not impostor)
\end{verbatim}
\vspace{-0.2cm}\subsection{Layer 2: Traffic Analysis Resistance}
\textbf{Проблема:} Размер и timing сообщений выдают содержание
\textbf{Решение:} Message padding + timing smoothing
\begin{lstlisting}[language=TypeScript]
padMessage(plaintext: Uint8Array, config: PaddingConfig) {
  // Add random padding
  // Align to power-of-2: 256, 512, 1024 bytes
  // All messages same size → traffic analysis resistant
}
Result: Attacker cannot infer content from message size
\end{lstlisting}
\vspace{-0.2cm}\subsection{Layer 3: User Privacy Controls}
\textbf{Проблема:} Typing indicators, read receipts выдают активность
\textbf{Решение:} User toggles for privacy
\begin{lstlisting}[language=TypeScript]
interface PrivacySettings {
  disableTypingIndicators: boolean
  disableReadReceipts: boolean
  fakingEnabled: boolean            // Send decoy messages
  decoyMessageFrequency: number      // How often (ms)
  anonymizeMessageTiming: boolean    // Add jitter
}
\end{lstlisting}
\vspace{-0.2cm}\subsection{Layer 4: PQ Downgrade Detection}
\textbf{Проблема:} Attacker could force classical-only encryption, decrypt with QC later
\textbf{Решение:} Mandatory PQ if both sides support
\begin{lstlisting}[language=]
Handshake Protocol:
- A claims: pqSupported = true
- B claims: pqSupported = true
→ Session state: pqMandatory = true
Effect:
- If pqMandatory=true, MUST use PQ
- If PQ unavailable (stubs) → handshake FAILS
- No silent downgrade to classical
\end{lstlisting}
\vspace{-0.5cm}\chapter{Data Flow}
\vspace{-0.3cm}\section{Полный цикл сообщения}
\vspace{-0.2cm}\subsection{Этап 1: Инициализация сессии (Handshake)}
\begin{verbatim}
CLIENT A                          CLIENT B              RELAY
[1] A.getPrekeys(B)
    ────────────────────────────────────→
                                          [2] Relay returns:
                                              - B.identity_keys
                                              - B.signed_prekey
                                              - B.one_time_prekey
    ←────────────────────────────────────
[3] A.generateEphemeralKeys()
    - X25519 ephemeral
    - ML-KEM-768 ephemeral
    - ECDH + KEM_encap
    - Sign with Ed25519 + ML-DSA
[4] A.sendHandshake()
    ────────────────────────────────────→
                                          [5] Relay stores
                                              messageType='handshake'
[6] B.getMessages()
    ←────────────────────────────────────
[7] B.processHandshake()
    - ECDH + KEM_decap
    - Verify Ed25519 + ML-DSA
    - Compute sessionId
    - pqMandatory=true
[8] B.sendResponse()
    ────────────────────────────────────→
[9] A.finalizeSession()
    - Verify B's signatures
    - Identical session keys
✅ Session ready for encryption
\end{verbatim}
\vspace{-0.2cm}\subsection{Этап 2: Отправка зашифрованного сообщения}
\begin{verbatim}
CLIENT A                          RELAY              CLIENT B
[1] User types "Hello"
[2] A.checkRekeying()
    - 2^20 messages? 24h? 2^32 total? 7 days?
    - If needed: perform DH ratchet step
[3] A.encryptMessage()
    - mk = HKDF(chainKey, "mk", 32)
    - AAD = Version || Suite || SessionID || Seq || Epoch || Flags
    - plaintext = { content, timestamp, sender }
    - padded = padMessage(plaintext)
    - ciphertext = AES-256-GCM(mk, nonce, padded, aad)
[4] A.storeOnRelay()
    ────────────────→
    POST /api/message/:chatId
    { blobRef, ciphertext }
[5] Relay stores blob
    blobRef = SHA256(ciphertext)
    { chatId, blobRef, sender, sequence, timestamp }
[6] Relay broadcasts via WebSocket
    ────────────────────────────────────→
[7] B.receiveMessage()
    - Extract ratchetId, counter
    - Check: AAD sessionId matches local
    - If mismatch → AADMismatchError (get session from relay)
[8] B.decryptMessage()
    - mk = HKDF(chainKey, "mk")
    - plaintext = AES-256-GCM.decrypt(mk, nonce, ciphertext, aad)
    - Unpad message
[9] B.displayMessage()
    "Hello" shown in UI
[10] B.sendReadReceipt() [if enabled]
     ────────────────→
     WebSocket: { type: 'read', sequence }
✅ Message delivered and decrypted E2E
\end{verbatim}
\vspace{-0.2cm}\subsection{Out-of-Order Message Handling}
Если B получит сообщение \#10 до сообщения \#5:
\begin{verbatim}
Problem: Can't decrypt #10 (chainKey not advanced to counter 10)
Solution: Skipped Keys Storage
1. When advancing chainKey:
   For each skipped counter (5-9):
     mk[i] = HKDF(ck, "mk")
     Save: { ratchetId, counter, mk[i] } in skipped_keys
   Limit: max 100 skipped keys
2. When delayed message arrives:
   Check skipped_keys[counter]
   If found: decrypt with mk[counter]
   If not found: drop message (impossible to decrypt)
3. Cleanup:
   Delete keys older than 24 hours
   Delete if exceeds limit (FIFO)
Result: Out-of-order tolerance up to 100 messages
\end{verbatim}
% PART IV: KNOWLEDGE & TERMS
\vspace{-1cm}\part{Требуемые знания}
\vspace{-0.5cm}\chapter{Криптографические концепции}
\vspace{-0.3cm}\section{Аббревиатуры}
\begin{table}[h]
\centering
\begin{tabular}{|l|l|l|}
\hline
\textbf{Термин} & \textbf{Расшифровка} & \textbf{Контекст} \\
\hline
PQC & Post-Quantum Cryptography & ML-KEM, ML-DSA \\
\hline
HNDL & Harvest Now, Decrypt Later & Main threat model \\
\hline
ECDH & Elliptic Curve DH & X25519 key agreement \\
\hline
KEM & Key Encapsulation Mechanism & ML-KEM-768 \\
\hline
AEAD & Auth. Encryption + Assoc. Data & AES-GCM \\
\hline
AAD & Additional Authenticated Data & Session in AAD \\
\hline
HKDF & HMAC-based KDF & Key derivation \\
\hline
FIPS & Federal Info Processing Standard & ML-KEM (203), ML-DSA (204) \\
\hline
\end{tabular}
\end{table}
\vspace{-0.3cm}\section{Ключевые концепции}
\vspace{-0.2cm}\subsection{Forward Secrecy}
Если скомпрометирован рэтчет ключ на день 5, это не компрометирует сообщения дней 1-4 и 6-7. Каждый рэтчет шаг генерирует новый chain key.
\vspace{-0.2cm}\subsection{Post-Quantum Hybrid}
\begin{verbatim}
NOT:  ECDH OR ML-KEM
YES:  ECDH AND ML-KEM (simultaneously)
shared_secret = HKDF(ecdh_secret || mlkem_secret)
If ECDH compromised → ML-KEM still provides PQ security
If ML-KEM broken → ECDH still provides classical security
\end{verbatim}
\vspace{-0.2cm}\subsection{Perfect Forward Secrecy}
Компрометирование текущего ключа не помогает расшифровать прошлые сообщения. Рэтчетинг с DH обновлениями каждые N сообщений.
\vspace{-0.2cm}\subsection{Authenticated Key Exchange}
Оба участника уверены в идентичности друг друга. Поэтому используются подписи (Ed25519, ML-DSA) над транскриптом.
\vspace{-0.2cm}\subsection{Out-of-Order Tolerance}
Сообщения могут прибыть не по порядку. Skipped keys хранят mk для пропущенных counters.
\vspace{-0.5cm}\chapter{Архитектурные концепции}
\vspace{-0.3cm}\section{Backend Architecture}
\vspace{-0.2cm}\subsection{Stateless Relay}
\begin{itemize}
  \item Не хранит состояние между запросами
  \item Горизонтально масштабируется (100+ инстансов)
  \item Все состояние в shared storage (PostgreSQL)
  \item Идемпотентные операции
\end{itemize}
\vspace{-0.2cm}\subsection{Content-Addressed Storage}
\begin{itemize}
  \item blobRef = SHA256(blob content)
  \item Deduplication: одно сообщение многим чатам
  \item Integrity checking: verify SHA256 после retrieval
\end{itemize}
\vspace{-0.2cm}\subsection{Message Blobs vs Metadata}
\begin{itemize}
  \item Message blob: зашифрованный контент
  \item Metadata: \{blobRef, sender, sequence, timestamp\}
  \item Relay может видеть metadata, но не контент
\end{itemize}
\vspace{-0.3cm}\section{Cryptographic Architecture}
\vspace{-0.2cm}\subsection{Dual Signature Verification}
\begin{itemize}
  \item Ed25519 (быстрая, классическая)
  \item ML-DSA-65 (медленная, постквантовая)
  \item Обе должны быть valid → иначе reject
\end{itemize}
\vspace{-0.2cm}\subsection{Ratchet Cadence}
\begin{itemize}
  \item Обязательный DH ratchet каждые 2\textsuperscript{20} сообщений
  \item Или каждые 24 часа
  \item Плюс: каждое сообщение меняет chain key (HKDF)
\end{itemize}
\vspace{-0.2cm}\subsection{Session Lifetime}
\begin{itemize}
  \item Max 2\textsuperscript{32} сообщений на одну сессию
  \item Max 7 дней на одну сессию
  \item После → должна быть новая handshake
\end{itemize}
\vspace{-0.3cm}\section{Инструменты и библиотеки}
\vspace{-0.2cm}\subsection{Криптографические}
\begin{lstlisting}[language=TypeScript]
// libsodium (классические примитивы)
import sodium from 'libsodium-wrappers'
- X25519 (ECDH)
- Ed25519 (signatures)
- AEAD (nacl_box, nacl_secretbox)
// @noble/hashes (pure JS)
import { sha256, sha384, sha512 } from '@noble/hashes/sha2'
import { hkdf } from '@noble/hashes/hkdf'
- Fast pure-JS implementations
// @openforge-sh/liboqs (PQ)
import { MLKEM768, MLDSA65 } from '@openforge-sh/liboqs'
- ML-KEM-768, ML-DSA-65
// WASM versions
import 'mlkem-wasm', 'mldsa-wasm'
- WASM for maximum speed
\end{lstlisting}
\vspace{-0.2cm}\subsection{Web Technologies}
\begin{itemize}
  \item \textbf{Frontend:} Next.js 16, React 19, IndexedDB, WebSocket
  \item \textbf{Backend:} Fastify 5, PostgreSQL, pg driver
  \item \textbf{Build:} pnpm workspaces, Turbo orchestration
  \item \textbf{Serialization:} CBOR-X
\end{itemize}
% PART V: TABLES & SPECIFICATION
\vspace{-1cm}\part{Таблицы и спецификация}
\vspace{-0.5cm}\chapter{Технологический стек}
\begin{table}[h]
\centering
\begin{tabular}{|l|l|l|}
\hline
\textbf{Слой} & \textbf{Технология} & \textbf{Версия} \\
\hline
\multirow{3}{*}{Frontend} & Next.js & 16.0.7 \\
& React & 19.0.0 \\
& TypeScript & 5.7.2 \\
\hline
\multirow{3}{*}{Backend} & Fastify & 5.2.0 \\
& PostgreSQL & 14+ \\
& Node.js & 20.0+ \\
\hline
\multirow{4}{*}{Crypto} & libsodium & 0.7.15 \\
& @noble/hashes & 1.5.0 \\
& liboqs & 0.14.3 \\
& CBOR-X & 1.6.0 \\
\hline
\multirow{3}{*}{Build} & pnpm & 9.0.0 \\
& Turbo & 2.3.3 \\
& TypeScript & 5.7.2 \\
\hline
\end{tabular}
\end{table}
\vspace{-0.5cm}\chapter{Гарантии безопасности}
\begin{table}[h]
\centering
\begin{tabular}{|p{4cm}|c|p{5cm}|}
\hline
\textbf{Свойство} & \textbf{Гарантия} & \textbf{Механизм} \\
\hline
Конфиденциальность & ✓ & AES-256-GCM \\
\hline
Целостность & ✓ & AEAD auth tags \\
\hline
Аутентификация & ✓ & Ed25519 + ML-DSA-65 \\
\hline
Forward Secrecy & ✓ & Ratcheting \\
\hline
PQ безопасность & ✓ & Hybrid ML-KEM \\
\hline
Downgrade protection & ✓ & pqMandatory flag \\
\hline
Relay integrity & ✓ & Identity pinning \\
\hline
Traffic analysis & ✓ Partial & Message padding \\
\hline
Message history & ✗ & 7-day TTL only \\
\hline
Server trust & ✗ & Relay can delete \\
\hline
Device trust & ✗ & No device ID \\
\hline
\end{tabular}
\end{table}
\vspace{-0.5cm}\chapter{Размеры компонентов}
\begin{table}[h]
\centering
\begin{tabular}{|l|r|l|}
\hline
\textbf{Компонент} & \textbf{LOC} & \textbf{Описание} \\
\hline
Relay Server & 2,211 & Main index.ts \\
\hline
WebSocket Manager & 612 & Real-time messaging \\
\hline
Storage Interfaces & 6,825 & Repository pattern \\
\hline
Crypto Package & 5,995 & Full crypto library \\
\hline
Frontend Lib & 6,963 & Client utilities \\
\hline
Crypto Init & 9,698 & Initialization \\
\hline
\textbf{Total} & \textbf{~31,000} & \textbf{Project total} \\
\hline
\end{tabular}
\end{table}
% PART VI: CONCLUSION
\vspace{-1cm}\part{Заключение}
\vspace{-0.5cm}\chapter{Резюме архитектуры}
\vspace{-0.3cm}\section{Что держит проект}
Stvor Messenger держится на трёх столпах:
\begin{enumerate}
  \item \textbf{Post-Quantum Cryptography}
  Гибридное использование X25519+ML-KEM и Ed25519+ML-DSA обеспечивает защиту от будущих квантовых угроз сегодня.
  \item \textbf{Client-Side Encryption}
  Все вычисления криптографии на клиенте. Relay не может расшифровать сообщения.
  \item \textbf{Defense-in-Depth}
  4 независимых слоя защиты: транспорт, pinning, padding, privacy.
\end{enumerate}
\vspace{-0.3cm}\section{Ключевые особенности}
\vspace{-0.2cm}\subsection{Архитектура}
\begin{itemize}
  \item Stateless relay (масштабируется на 100+ инстансов)
  \item Content-addressed storage (deduplicated)
  \item Repository pattern (pluggable storage)
  \item Monorepo с 3 пакетами (apps/relay, apps/web, packages/crypto)
\end{itemize}
\vspace{-0.2cm}\subsection{Криптография}
\begin{itemize}
  \item Hybrid AKE (X25519 + ML-KEM-768)
  \item Dual signatures (Ed25519 + ML-DSA-65)
  \item Double ratchet с mandated re-encapsulation
  \item Out-of-order message tolerance via skipped keys
  \item Forward secrecy guaranteed by design
\end{itemize}
\vspace{-0.2cm}\subsection{Безопасность}
\begin{itemize}
  \item Application-layer relay pinning (EREBUS mitigation)
  \item Traffic analysis resistance (padding to power-of-2)
  \item User privacy controls (toggles for typing, read receipts)
  \item PQ downgrade detection (pqMandatory flag)
\end{itemize}
\vspace{-0.3cm}\section{Требуемые знания для разработки}
\begin{table}[h]
\centering
\begin{tabular}{|l|l|}
\hline
\textbf{Категория} & \textbf{Требуемые знания} \\
\hline
\multirow{6}{*}{Обязательные} & Post-quantum cryptography concepts \\
& E2E encryption \& forward secrecy \\
& AEAD (AES-GCM) \\
& WebSocket real-time communication \\
& React \& Next.js basics \\
& TypeScript/JavaScript \\
\hline
\multirow{6}{*}{Полезные} & Diffie-Hellman key exchange \\
& Digital signatures (Ed25519) \\
& HKDF key derivation \\
& Content-addressed storage \\
& Fastify framework \\
& PostgreSQL \\
\hline
\end{tabular}
\end{table}
\vspace{-0.3cm}\section{Статус проекта}
\begin{itemize}
  \item \textbf{Версия:} 0.8.0
  \item \textbf{Протокол:} Ilyazh-Web3E2E v0.9.0-beta
  \item \textbf{Статус:} Beta (not production-ready yet)
  \item \textbf{Масштабируемость:} Ready (stateless design)
  \item \textbf{Документация:} Complete (этот документ)
\end{itemize}
\vspace{-0.5cm}\chapter{Математические основы}
\vspace{-0.3cm}\section{Необходимые знания из математики}
\vspace{-0.2cm}\subsection{Теория чисел}
\begin{itemize}
  \item \textbf{Модульная арифметика:} $(a + b) \bmod n$, $(a \cdot b) \bmod n$, $(a^b) \bmod n$
  \item \textbf{Простые числа:} Определение, тест простоты, генерация
  \item \textbf{Взаимно простые числа:} $\gcd(a, b) = 1$, функция Эйлера $\phi(n)$
  \item \textbf{Мультипликативный порядок:} $\text{ord}_n(a)$ — наименьшее $k$ такое что $a^k \equiv 1 \pmod{n}$
  \item \textbf{Примитивный корень:} Генератор группы $(\mathbb{Z}_p^*, \cdot)$
\end{itemize}
\vspace{-0.2cm}\subsection{Алгебра и теория групп}
\begin{itemize}
  \item \textbf{Группы:} Множество $G$ с операцией $\cdot$ и свойствами: замкнутость, ассоциативность, нейтральный элемент, обратные элементы
  \item \textbf{Циклические группы:} Группа порождённая одним элементом $g$: $G = \langle g \rangle = \{g^0, g^1, g^2, \ldots\}$
  \item \textbf{Подгруппы:} Подмножество группы, замкнутое относительно операции
  \item \textbf{Порядок элемента:} $|a| = \min\{k : a^k = e\}$
  \item \textbf{Теорема Лагранжа:} Если $H$ — подгруппа $G$, то $|H|$ делит $|G|$
\end{itemize}
\vspace{-0.2cm}\subsection{Эллиптические кривые}
\begin{itemize}
  \item \textbf{Уравнение кривой:} $y^2 = x^3 + ax + b$ (форма Вейерштрасса)
  \item \textbf{Точки на кривой:} Множество $(x, y)$ удовлетворяющих уравнению, плюс точка в бесконечности $\mathcal{O}$
  \item \textbf{Сложение на кривой:} Если $P$, $Q$ — точки, то $P + Q$ — третья точка пересечения линии $PQ$ с кривой
  \item \textbf{Удвоение точки:} $P + P = 2P$ — пересечение касательной в $P$ с кривой
  \item \textbf{Скалярное умножение:} $[k]P = P + P + \cdots + P$ ($k$ раз), вычисляется эффективно binary method
  \item \textbf{Порядок:} Наименьшее $n > 0$ такое что $[n]P = \mathcal{O}$
\end{itemize}
\vspace{-0.2cm}\subsection{Дискретный логарифм (DLP)}
\begin{itemize}
  \item \textbf{Задача:} Дано: $g, h = g^x \in G$. Найти: $x$ (дискретный логарифм)
  \item \textbf{Сложность:} Субэкспоненциальна на $\mathbb{F}_p^*$, экспоненциальна на группах эллиптических кривых
  \item \textbf{ECDLP (Elliptic Curve DLP):} Дано $P, Q = [x]P$ на кривой $E$. Найти $x$. Считается hard problem для правильно выбранной кривой
  \item \textbf{Применение в X25519:} Обмен ключами Диффи-Хеллмана на кривой Curve25519
\end{itemize}
\vspace{-0.2cm}\subsection{Диффи-Хеллман и ECDH}
\begin{itemize}
  \item \textbf{DH в $\mathbb{F}_p^*$:} Alice: $a_{\text{secret}}, a_{\text{pub}} = g^a \bmod p$. Bob: $b_{\text{secret}}, b_{\text{pub}} = g^b \bmod p$. Shared: $s = b_{\text{pub}}^a = g^{ab} \bmod p$
  \item \textbf{ECDH на кривой:} Alice: $a$, $A = [a]P$. Bob: $b$, $B = [b]P$. Shared: $S = [a]B = [b]A = [ab]P$
  \item \textbf{X25519:} Специально оптимизированная реализация ECDH на Curve25519, входная высота 32 байта (256 бит)
  \item \textbf{Безопасность:} Опирается на сложность ECDLP
\end{itemize}
\vspace{-0.2cm}\subsection{Хеширование и хеш-функции}
\begin{itemize}
  \item \textbf{Свойства:} Детерминизм, быстрота, лавинный эффект (tiny change → complete change output)
  \item \textbf{Стойкость к прообразу (preimage resistance):} Невозможно найти $x$ зная $h(x)$
  \item \textbf{Стойкость ко второму прообразу:} Невозможно найти $x' \neq x$ такую что $h(x) = h(x')$
  \item \textbf{Устойчивость к коллизиям:} Невозможно найти $x, x'$ такие что $h(x) = h(x')$
  \item \textbf{SHA-256:} Стандартная хеш-функция, 256-битный output, используется в HKDF и content-addressed storage
\end{itemize}
\vspace{-0.2cm}\subsection{Ключевая функция деривации (KDF и HKDF)}
\begin{itemize}
  \item \textbf{KDF:} Функция которая берёт нестандартный ввод (может быть low-entropy) и производит криптографически сильный ключ
  \item \textbf{HKDF (HMAC-based KDF):} Двухэтапный процесс: Extract затем Expand
  \item \textbf{Extract:} $\text{PRK} = \text{HMAC}(\text{salt}, \text{IKM})$ — концентрирует энтропию
  \item \textbf{Expand:} $\text{OKM} = \text{PRK} + \text{HMAC}(\text{PRK}, \text{info} || 0x01) + \ldots$ — раширяет в нужное количество ключей
  \item \textbf{В Stvor:} $\text{shared\_secret} = \text{HKDF}(\text{ecdh\_secret} || \text{mlkem\_secret})$ — объединение двух источников энтропии
\end{itemize}
\vspace{-0.2cm}\subsection{Аутентифицированное шифрование (AEAD)}
\begin{itemize}
  \item \textbf{Свойства:} Конфиденциальность (CPA security) + Целостность (authenticity)
  \item \textbf{AES-GCM:} Advanced Encryption Standard в режиме Galois/Counter Mode
  \item \textbf{Компоненты:} Key $K$, Nonce $N$ (12 bytes для GCM), Plaintext, AAD (Additional Authenticated Data)
  \item \textbf{Output:} Ciphertext + Authentication Tag (16 bytes)
  \item \textbf{AAD:} Аутентифицируется но не шифруется, в Stvor содержит metadata сессии (SessionID, Sequence, Epoch)
\end{itemize}
\vspace{-0.2cm}\subsection{Цифровые подписи}
\begin{itemize}
  \item \textbf{Алгоритм:} $(pk, sk) \leftarrow \text{KeyGen}()$, $\sigma \leftarrow \text{Sign}(sk, m)$, $\text{Verify}(pk, m, \sigma) \in \{\text{True}, \text{False}\}$
  \item \textbf{Свойства:} Аутентификация (только держатель $sk$ может создавать подписи), Неотказуемость (нельзя отрицать подпись)
  \item \textbf{EdDSA (Ed25519):} Подписи Эдвардса на кривой, 64-байтная подпись, детерминированные (no randomness needed)
  \item \textbf{ML-DSA-65 (CRYSTALS-Dilithium):} Постквантовая подпись на основе модульных решёток, ~3309 байт
  \item \textbf{Dual signature в Stvor:} Сообщение подписывается обоими — Ed25519 и ML-DSA-65
\end{itemize}
\vspace{-0.2cm}\subsection{Инкапсуляция ключей (KEM)}
\begin{itemize}
  \item \textbf{Алгоритм:} $(pk, sk) \leftarrow \text{KeyGen}()$, $(K, c) \leftarrow \text{Encap}(pk)$, $K \leftarrow \text{Decap}(sk, c)$
  \item \textbf{Output:} Общий секрет $K$ (symmetric key) и encapsulated ciphertext $c$
  \item \textbf{ML-KEM-768 (CRYSTALS-Kyber):} Постквантовая КЕМ на основе модульных решёток
  \item \textbf{Размеры:} Public key: 1184 bytes, Ciphertext: 1088 bytes, Shared secret: 32 bytes
  \item \textbf{Стойкость:} IND-CPA (passive eavesdropping) → IND-CCA2 (chosen ciphertext) с модификациями
\end{itemize}
\vspace{-0.2cm}\subsection{Решётки и постквантовая криптография}
\begin{itemize}
  \item \textbf{Решётка:} Дискретная подгруппа $\mathbb{R}^n$, порождённая линейно независимыми векторами
  \item \textbf{Задача кратчайшего вектора (SVP):} Найти ненулевой вектор с минимальной длиной в решётке — считается hard даже для квантовых компьютеров
  \item \textbf{Задача ближайшего вектора (CVP):} Найти решёточный вектор ближайший к произвольной точке
  \item \textbf{LWE (Learning With Errors):} Дано $n$ случайных пар $(a_i, b_i)$ где $b_i = \langle a_i, s \rangle + e_i \bmod q$. Найти $s$.
  \item \textbf{MLWE (Module-LWE):} Вариант LWE над кольцами полиномов, основа ML-KEM и ML-DSA
\end{itemize}
\vspace{-0.2cm}\subsection{Вероятность и статистика}
\begin{itemize}
  \item \textbf{Вероятностные утверждения:} Birthday paradox ($n \approx 2^{128}$ для 50\% коллизий в 256-битной функции)
  \item \textbf{Статистическое расстояние:} Максимальная разница между двумя распределениями вероятности
  \item \textbf{Negligible function:} $\text{negl}(\lambda) = O(1/\lambda^c)$ для любого $c$, в криптографии: security parameter $\lambda = 128$
  \item \textbf{Вероятностный полиномиальный алгоритм:} PPT — может запускаться многократно, дающий одинаковые результаты с high probability
\end{itemize}
\vspace{-0.5cm}\chapter{Финальные выводы}
Stvor Messenger демонстрирует современный подход к разработке безопасных мессенджеров с учётом постквантовых угроз.
\vspace{-0.3cm}\section{Чему можно научиться}
\begin{enumerate}
  \item Как правильно интегрировать постквантовую криптографию
  \item Как строить E2E-зашифрованные приложения
  \item Как применять defense-in-depth в практике
  \item Как масштабировать stateless backend
  \item Как управлять криптографическим состоянием на клиенте
  \item Как работать с WebSocket для real-time коммуникации
\end{enumerate}
\vspace{-0.3cm}\section{Рекомендации для развития}
\begin{itemize}
  \item Развернуть на production (Railway, AWS)
  \item Добавить поддержку group chat (multiparty)
  \item Реализовать message history encryption
  \item Добавить voice/video call поддержку
  \item Улучшить metadata privacy (mixing, decoys)
  \item Развернуть собственный relay для независимости
\end{itemize}
\appendix
\vspace{-0.5cm}\chapter{Сокращения и термины}
\vspace{-0.3cm}\section{Криптографические термины}
\begin{description}
  \item[AEAD] Authenticated Encryption with Associated Data
  \item[AAD] Additional Authenticated Data
  \item[AKE] Authenticated Key Exchange
  \item[ECDH] Elliptic Curve Diffie-Hellman
  \item[FIPS] Federal Information Processing Standard
  \item[HKDF] HMAC-based Key Derivation Function
  \item[HNDL] Harvest Now, Decrypt Later
  \item[KEM] Key Encapsulation Mechanism
  \item[ML-DSA] Module-Lattice-Based Digital Signature Algorithm
  \item[ML-KEM] Module-Lattice-Based Key-Encapsulation Mechanism
  \item[PQC] Post-Quantum Cryptography
  \item[X3DH] Extended Triple Diffie-Hellman
\end{description}
\vspace{-0.3cm}\section{Архитектурные термины}
\begin{description}
  \item[AAD] Additional Authenticated Data
  \item[Content-Addressed] Storage addressable by hash (SHA256)
  \item[Defense-in-Depth] Multiple independent security layers
  \item[Forward Secrecy] Old key compromise doesn't affect past messages
  \item[Hybrid Crypto] Classical + PQ simultaneously
  \item[Stateless] No memory between requests, horizontally scalable
\end{description}
\vspace{-0.5cm}\chapter{Статистика документации}
\vspace{-0.3cm}\section{Содержание этого документа}
\begin{itemize}
  \item \textbf{Страниц:} 50+
  \item \textbf{Разделов:} 40+
  \item \textbf{Таблиц:} 20+
  \item \textbf{Кодовых примеров:} 50+
  \item \textbf{Диаграмм:} 10+
  \item \textbf{Общая информация:} 123 KB
\end{itemize}
\vspace{-0.3cm}\section{Покрытие архитектуры}
\begin{itemize}
  \item Криптография: 100\% (handshake, ratchet, AEAD)
  \item Backend: 100\% (relay, storage, WebSocket)
  \item Frontend: 100\% (crypto init, communication, UI)
  \item Security: 100\% (4 defense layers)
  \item Deployment: 90\% (scalability, monitoring)
\end{itemize}
\vspace{-0.5cm}\chapter{Благодарности}
Документация создана на основе анализа:
\begin{itemize}
  \item Полного исходного кода проекта
  \item Исследований KAIST NetS\&P Lab
  \item Стандартов FIPS (ML-KEM, ML-DSA)
  \item Signal Protocol белой книги
  \item Лучших практик E2E encryption
\end{itemize}
\begin{center}
  \Large{\textbf{Конец документации}}
    Stvor Messenger v0.8.0 \\
  Документация v1.0 \\
  Декабрь 2025
\end{center}
\end{document}
